\documentclass{article}
\usepackage[utf8]{inputenc}
\usepackage[T1]{fontenc}
\usepackage[brazil]{babel}
\usepackage{graphicx}
\usepackage{booktabs}
\usepackage{algorithm}
\usepackage{algpseudocode}

\title{Seleção de Serviços sob Incerteza --- Mergulho Inviável vs. Híbrido (Best-Fit + Oscilação)}
\author{Ramon Barros de Lima}
\date{Julho 2026}

\begin{document}

\maketitle

\section{O problema}

Devemos alocar 100 tarefas em 100 serviços gastando o \textbf{mínimo}, respeitando três
restrições que não podem ser quebradas: (i) \textbf{custo} total; (ii) \textbf{capacidade}
de recurso de cada serviço ($V_{res}$); e (iii) \textbf{SLA}, cuja probabilidade de violação
não pode passar de $P_{max}$. O algoritmo base é um \textit{Iterated Local Search} (ILS):
constrói uma solução inicial e a refina alternando \textbf{perturbação} e \textbf{busca local}.

Comparamos duas estratégias para escapar de ótimos locais nas instâncias mais difíceis:
o \textbf{Mergulho Inviável} (Caso A) e o \textbf{Híbrido best-fit + oscilação} (Caso 3).

\section{Caso A --- Mergulho Inviável (IGrAl)}

Quando a busca fica presa, esta estratégia \textbf{mergulha em soluções inviáveis} (violando o
SLA / $P_{max}$) na esperança de escapar para outra região. A solução inviável é penalizada
por $\lambda = \text{custo}(s^*)/P_{max}$; se o custo penalizado for promissor, roda-se a busca
local a partir dela; senão, reinicia-se com um construtivo adaptativo.

\begin{algorithm}[H]
\caption{ILS com Mergulho em Soluções Inviáveis (Caso A)}
\begin{algorithmic}[1]
\Require instância, $\alpha$, $IT\_MAX$
\Ensure melhor solução viável $s^*$

\State $s \leftarrow \text{SoluçãoGulosa}(\text{instância}, \alpha)$
\State $s \leftarrow \text{BuscaLocal}(s)$
\State $s^* \leftarrow s$
\State $semMelhora \leftarrow 0$;\quad $reiniciosSemMelhora \leftarrow 0$

\For{$i = 1$ \textbf{até} $IT\_MAX$}
    \State $s \leftarrow \text{Perturbação}(s,\ \textsc{Move})$
    \State $s \leftarrow \text{BuscaLocal}(s)$
    \If{$\text{custo}(s) < \text{custo}(s^*)$}
        \State $s^* \leftarrow s$;\quad $semMelhora \leftarrow 0$;\quad $reiniciosSemMelhora \leftarrow 0$
    \Else
        \State $semMelhora \leftarrow semMelhora + 1$
    \EndIf

    \If{$semMelhora > IT\_MAX / 10$}
        \State $semMelhora \leftarrow 0$
        \State $dive \leftarrow \text{Perturbação}(s,\ \textsc{Infeasible})$
        \Comment{ignora $P_{max}$}
        \State $excesso \leftarrow \text{ComputaExcessoViolação}(dive)$
        \State $\lambda \leftarrow \text{custo}(s^*)\ /\ P_{max}$
        \State $custoP \leftarrow \text{custo}(dive) + \lambda \times excesso$
        \State $w \leftarrow \min(0.5,\ reiniciosSemMelhora\ /\ 5)$

        \If{$custoP < \text{custo}(s^*)$}
            \Comment{dive promissor (raro na prática)}
            \State $explorado \leftarrow \text{BuscaLocal}(dive)$
            \If{$\text{viável}(explorado)$}
                \If{$\text{custo}(explorado) < \text{custo}(s^*)$}
                    \State $s^* \leftarrow explorado$;\quad $reiniciosSemMelhora \leftarrow 0$
                \EndIf
                \State $s \leftarrow explorado$
            \Else
                \State $s \leftarrow \text{ConstrutivoAdaptativo}(\text{instância}, \alpha, w)$
                \State $reiniciosSemMelhora \leftarrow reiniciosSemMelhora + 1$
            \EndIf
        \Else
            \State $s \leftarrow \text{ConstrutivoAdaptativo}(\text{instância}, \alpha, w)$
            \State $reiniciosSemMelhora \leftarrow reiniciosSemMelhora + 1$
        \EndIf
    \EndIf
\EndFor
\State \Return $s^*$
\end{algorithmic}
\end{algorithm}

\noindent\textbf{Limitação medida:} o dive relaxa o \textbf{SLA}, mas o gargalo das instâncias
difíceis é a \textbf{capacidade}. Além disso, $\lambda$ é alto demais e o custo penalizado quase
nunca fica abaixo de $\text{custo}(s^*)$ --- o ramo do ``dive promissor'' praticamente não dispara.
Na prática, degenera em reinícios.

\section{Caso 3 --- Híbrido (Best-Fit + Oscilação Estratégica)}

O híbrido substitui \emph{duas} peças do ILS: o \textbf{construtivo} (Proposta 1) e a
\textbf{busca local} (Proposta 2).

\paragraph{Proposta 1 --- Construtivo Best-Fit.} Aloca as tarefas mais \textbf{pesadas} primeiro
(as ``pedras grandes'') e, para cada uma, escolhe o serviço \textbf{mais barato que ainda mantém
a solução viável} (\textit{cheapest-feasible}) --- checando as restrições \emph{antes} de fixar.
Parte de uma estrutura bem ``empacotada''.

\paragraph{Proposta 2 --- Oscilação Estratégica.} Relaxa a \textbf{capacidade} para restrição
\emph{suave}, otimizando $f = \text{custo} + \lambda \cdot \text{sobrecarga}$. O multiplicador
$\lambda$ \textbf{oscila}: diminui quando a solução é viável (deixa cruzar para regiões
inviáveis-mas-baratas) e aumenta quando é inviável (empurra de volta). Isso permite o movimento
que a busca estritamente viável nunca faz --- sobrecarregar um serviço temporariamente para
reequilibrar e eliminar outro depois. $S_{max}$ e SLA permanecem restrições duras.

\begin{algorithm}[H]
\caption{ILS Híbrido (Caso 3)}
\begin{algorithmic}[1]
\Require instância, $\alpha$, $IT\_MAX$
\Ensure melhor solução viável $s^*$
\State $s \leftarrow \textsc{ConstrutivoBestFit}(\text{instância}, 0)$
\State $s \leftarrow \textsc{BuscaLocalHíbrida}(s)$
\Comment{descida por custo $+$ oscilação}
\State $s^* \leftarrow s$
\For{$i = 1$ \textbf{até} $IT\_MAX$}
    \State $s \leftarrow \text{Perturbação}(s,\ \textsc{Move})$
    \State $s \leftarrow \textsc{BuscaLocalHíbrida}(s)$
    \If{$\text{custo}(s) < \text{custo}(s^*)$}
        \State $s^* \leftarrow s$
    \EndIf
\EndFor
\State \Return $s^*$
\end{algorithmic}
\end{algorithm}

\noindent onde $\textsc{BuscaLocalHíbrida}(s)$ aplica a descida por custo (aceita só moves viáveis
e mais baratos) seguida da \textsc{OscilaçãoEstratégica} (Algoritmo~4).

\begin{algorithm}[H]
\caption{ConstrutivoBestFit --- Proposta 1}
\begin{algorithmic}[1]
\Require instância, $\alpha$
\Ensure solução inicial $s$ (ou falha)
\State ordena tarefas por consumo de recurso \textbf{decrescente} $\rightarrow T$
\State $s \leftarrow \emptyset$
\For{cada tarefa $t \in T$}
    \State ordena serviços por custo \textbf{crescente} para $t$
    \State define a ordem de tentativa (RCL aleatória se $\alpha>0$; senão do mais barato ao mais caro)
    \State $colocada \leftarrow \textbf{falso}$
    \For{cada serviço $j$ na ordem de tentativa}
        \State atribui $t \rightarrow j$ em $s$
        \If{$\textsc{Viável}(s)$}
            \Comment{capacidade, $S_{max}$ e SLA}
            \State $colocada \leftarrow \textbf{verdadeiro}$;\quad \textbf{break}
        \Else
            \State desfaz a atribuição
        \EndIf
    \EndFor
    \If{\textbf{não} $colocada$}
        \State \Return falha
    \EndIf
\EndFor
\State \Return $s$
\end{algorithmic}
\end{algorithm}

\begin{algorithm}[H]
\caption{OscilaçãoEstratégica --- Proposta 2}
\begin{algorithmic}[1]
\Require solução viável $s$; número de rodadas $R$
\Ensure solução viável (melhorada ou a original)
\State $s_0 \leftarrow s$;\quad $melhor \leftarrow s$
\State $\lambda \leftarrow 2 \times \text{custoMédioMínimo}$
\For{$r = 1$ \textbf{até} $R$}
    \State $s \leftarrow \textsc{DescidaPenalizada}(s, \lambda)$
    \Comment{minimiza $f = \text{custo} + \lambda\cdot\text{sobrecarga}$ (Move)}
    \If{$\text{sobrecarga}(s) = 0$}
        \Comment{região viável}
        \If{$\text{custo}(s) < \text{custo}(melhor)$}
            \State $melhor \leftarrow s$
        \EndIf
        \State $\lambda \leftarrow \lambda \times 0.6$
        \Comment{relaxa}
    \Else
        \State $\lambda \leftarrow \lambda \times 2.0$
        \Comment{aperta}
    \EndIf
\EndFor
\State $s \leftarrow \textsc{DescidaPenalizada}(s, 8\lambda)$
\Comment{descida final garante viabilidade}
\If{$\text{sobrecarga}(s)=0$ \textbf{e} $\text{custo}(s) < \text{custo}(melhor)$}
    \State $melhor \leftarrow s$
\EndIf
\If{$\text{custo}(melhor) < \text{custo}(s_0)$}
    \State \Return $melhor$
\Else
    \State \Return $s_0$
\EndIf
\end{algorithmic}
\end{algorithm}

\noindent $\textsc{DescidaPenalizada}(s,\lambda)$ é uma descida \emph{best-improvement} por
movimentos \textsc{Move}: aplica repetidamente o move que mais reduz $f = \text{custo} +
\lambda\cdot\text{sobrecarga}$, aceitando sobrecarregar um serviço se o ganho líquido em $f$
compensar. $S_{max}$ e SLA são verificados e mantidos duros.

\section{A melhoria do Caso 3}

Comparação isolada no \emph{mesmo motor} ILS, variando apenas o mecanismo, nas 5 instâncias mais
difíceis (3 execuções cada):

\begin{table}[H]
\centering
\begin{tabular}{ccrr}
\toprule
Instância & Ótimo & Mergulho Inviável (GAP) & Híbrido (GAP) \\
\midrule
11  & 100 & 154 \ (54\%) & \textbf{108 \ (8\%)} \\
28  & 102 & 165 \ (62\%) & \textbf{123 \ (21\%)} \\
100 & 101 & 168 \ (66\%) & \textbf{131 \ (30\%)} \\
128 & 101 & 161 \ (59\%) & \textbf{116 \ (15\%)} \\
129 & 100 & 159 \ (59\%) & \textbf{108 \ (8\%)} \\
\midrule
\multicolumn{2}{c}{\textbf{GAP médio}} & \textbf{60,1\%} & \textbf{16,2\%} \\
\bottomrule
\end{tabular}
\caption{Mergulho inviável (viola SLA) vs. Híbrido (viola capacidade). Mesmo motor,
só a estratégia muda.}
\end{table}

O híbrido vence em \textbf{5 de 5} instâncias, reduzindo o GAP médio em cerca de \textbf{44 pontos
percentuais}: o \emph{pior} resultado do híbrido (30\%) ainda supera o \emph{melhor} do mergulho
inviável (54\%). Numa avaliação mais ampla, nas 94 instâncias, o híbrido leva o GAP médio de
$52{,}2\%$ (baseline) para $4{,}1\%$, atingindo o ótimo em $46/94$ instâncias.

\paragraph{Por que o Caso 3 é superior.} O mergulho inviável mexe na restrição \emph{errada}
(SLA) com um mecanismo que quase nunca dispara. O híbrido relaxa exatamente o \emph{gargalo}
(capacidade), com um $\lambda$ adaptativo que de fato aceita os movimentos de reequilíbrio,
partindo de uma construção best-fit bem estruturada.

\paragraph{Ressalva.} Mesmo o híbrido não atinge o ótimo perfeito nas mais apertadas (ex.\ a
instância 100 chega a $\approx 131$, com ótimo $101$): por ser uma busca por vizinhança, nunca
\emph{prova} otimalidade. Para essa garantia, seria necessário um método exato (ILP /
matheurística).

\end{document}